% paper/discussion.tex: the Discussion, pulled into main.tex with % paper/discussion.tex: the Discussion, pulled into main.tex with % paper/discussion.tex: the Discussion, pulled into main.tex with % paper/discussion.tex: the Discussion, pulled into main.tex with \input{discussion}.
% Edit freely; rebuild with: tectonic -X compile main.tex

The yes/no rule is overconfident mainly because it assumes that every lap before the leader responds earns the full fresh-tyre advantage, taken as a median over other cars rather than what this chaser gets on this compound. The first of those laps is the least reliable, since the chaser starts it on cold tyres while the leader can push on its in-lap, and we suspect that a good share of undercuts are decided there. Our data cannot settle whether that lap should count at all. With it, the rule's 50/50 point is a margin of 1.0~s; without it, the point moves to $-0.3$~s.

With two fitted numbers, the calibrated rule predicted as well as the full model with ten (Brier scores 0.211 and 0.218; Table~\ref{tab:rq2}), and the gap and the response delay alone did at least as well as either (0.200; Table~\ref{tab:rq5}). Since there are only 154 attempts and 77 successes, every extra coefficient carries noise, and in a new season that noise can cost more than the factor adds. The full model was the best method with 2023 held out but worse than the base rate with 2025 held out. Broadcasters could quote a probability instead of calling the outcome. Teams have more to go on (tyre temperatures, degradation models, the rival's radio and their own plans), and any of these might beat the gap alone.

Of the two numbers that matter, the chaser controls only the gap (Figure~\ref{fig:lookup}). From within a second, 79\% of undercuts worked, and the model gives better-than-even odds up to about 1.2~s even when the leader responds on the next lap. At 2~s the odds are better than even only if the leader waits three laps or more, so a chaser at the edge of the ``striking range'' is relying on a slow response. If leaders waited only when they were safe, undercuts against a slow response should rarely have worked, but 72--75\% of them did. Some leaders may be forced to wait by another car they are racing, a teammate already in the pit lane, traffic where they would rejoin or a plan for a longer stint, and others may simply misjudge the threat. Because the leader picks its delay after seeing the chaser stop, we cannot say that a late response causes a failed defence. Game-theoretic models of pit stops \citep{aguad2024pit} treat the undercut as a contest between two drivers' decisions, and our results are consistent with that.

About half of the difference between successful and failed undercuts came from the stops themselves (3.6 of 7.5~s; Figure~\ref{fig:mech}). When the leader's stop was more than 3~s slower, 15 of 17 undercuts worked. A team's pit-crew record did not predict its stops in a given attempt (correlation 0.07), so slow stops look like one-off events. An undercut that failed on a slow wheel change may therefore have been the right call when it was made, and should be judged that way. Race simulations should also treat stop time as random, as Monte Carlo simulations do \citep{heilmeier2020monte}, since fixed stop times would miss an effect of this size.

Faster cars won more undercuts (58\% against 44\%), but they also started closer, and once the gap and the response delay were in the model, pace was no longer associated with success (odds ratio 1.00 per 0.1~s a lap). The team effect that did show up was on the defending side (Figure~\ref{fig:teams}), though with fewer than 25 attempts per team and six tests added after the main analysis it is weak evidence. Undercuts against Aston Martin, Mercedes and Ferrari rarely worked, and those against Alpine, Racing Bulls and Sauber/Audi usually did. Some teams may be better at parts of a defence that timing data cannot see, such as warming their own new tyres on the out-lap. Front-runners also did worse as chasers. With the gap and the response delay held fixed, a chaser in the top six had lower odds of success (odds ratio 0.41, $p=0.06$, in groups chosen after seeing the data). The leading teams may watch each rival more closely and respond in ways the delay does not capture.

Circuits mattered mainly through the gap, with undercuts working more often where attempts started closer ($\rho=-0.54$) and no more often where tyre wear or the fresh-tyre gain was larger. A circuit's record also did not improve held-out predictions, which is unsurprising with fewer than ten attempts at most circuits.

In 2026 fewer undercuts worked (36\% against 50\%; Figure~\ref{fig:2026}) and leaders responded faster, with 14 of the 22 responses on the next lap; only 3 of those undercuts worked. Even with the delay and the other factors held fixed, 2026 attempts had lower odds (odds ratio 0.34), and the 2023--25 model expected 56\% success against the 36\% observed. The new tyres may take longer to warm up, which would hurt the chaser's out-lap, and teams may also have learned to respond sooner. The calibrated rule and the gap-and-delay model still beat the base rate in 2026 (Brier scores 0.201 and 0.212, against 0.250), so the arithmetic seems to have survived the new rules. We cannot yet say the odds have changed. The interval for the change in success rate runs from $-36$ to $+3$ points, which includes no change, and the 2026 indicator reaches only $p=0.049$ in a test we planned as exploratory.
.
% Edit freely; rebuild with: tectonic -X compile main.tex

The yes/no rule is overconfident mainly because it assumes that every lap before the leader responds earns the full fresh-tyre advantage, taken as a median over other cars rather than what this chaser gets on this compound. The first of those laps is the least reliable, since the chaser starts it on cold tyres while the leader can push on its in-lap, and we suspect that a good share of undercuts are decided there. Our data cannot settle whether that lap should count at all. With it, the rule's 50/50 point is a margin of 1.0~s; without it, the point moves to $-0.3$~s.

With two fitted numbers, the calibrated rule predicted as well as the full model with ten (Brier scores 0.211 and 0.218; Table~\ref{tab:rq2}), and the gap and the response delay alone did at least as well as either (0.200; Table~\ref{tab:rq5}). Since there are only 154 attempts and 77 successes, every extra coefficient carries noise, and in a new season that noise can cost more than the factor adds. The full model was the best method with 2023 held out but worse than the base rate with 2025 held out. Broadcasters could quote a probability instead of calling the outcome. Teams have more to go on (tyre temperatures, degradation models, the rival's radio and their own plans), and any of these might beat the gap alone.

Of the two numbers that matter, the chaser controls only the gap (Figure~\ref{fig:lookup}). From within a second, 79\% of undercuts worked, and the model gives better-than-even odds up to about 1.2~s even when the leader responds on the next lap. At 2~s the odds are better than even only if the leader waits three laps or more, so a chaser at the edge of the ``striking range'' is relying on a slow response. If leaders waited only when they were safe, undercuts against a slow response should rarely have worked, but 72--75\% of them did. Some leaders may be forced to wait by another car they are racing, a teammate already in the pit lane, traffic where they would rejoin or a plan for a longer stint, and others may simply misjudge the threat. Because the leader picks its delay after seeing the chaser stop, we cannot say that a late response causes a failed defence. Game-theoretic models of pit stops \citep{aguad2024pit} treat the undercut as a contest between two drivers' decisions, and our results are consistent with that.

About half of the difference between successful and failed undercuts came from the stops themselves (3.6 of 7.5~s; Figure~\ref{fig:mech}). When the leader's stop was more than 3~s slower, 15 of 17 undercuts worked. A team's pit-crew record did not predict its stops in a given attempt (correlation 0.07), so slow stops look like one-off events. An undercut that failed on a slow wheel change may therefore have been the right call when it was made, and should be judged that way. Race simulations should also treat stop time as random, as Monte Carlo simulations do \citep{heilmeier2020monte}, since fixed stop times would miss an effect of this size.

Faster cars won more undercuts (58\% against 44\%), but they also started closer, and once the gap and the response delay were in the model, pace was no longer associated with success (odds ratio 1.00 per 0.1~s a lap). The team effect that did show up was on the defending side (Figure~\ref{fig:teams}), though with fewer than 25 attempts per team and six tests added after the main analysis it is weak evidence. Undercuts against Aston Martin, Mercedes and Ferrari rarely worked, and those against Alpine, Racing Bulls and Sauber/Audi usually did. Some teams may be better at parts of a defence that timing data cannot see, such as warming their own new tyres on the out-lap. Front-runners also did worse as chasers. With the gap and the response delay held fixed, a chaser in the top six had lower odds of success (odds ratio 0.41, $p=0.06$, in groups chosen after seeing the data). The leading teams may watch each rival more closely and respond in ways the delay does not capture.

Circuits mattered mainly through the gap, with undercuts working more often where attempts started closer ($\rho=-0.54$) and no more often where tyre wear or the fresh-tyre gain was larger. A circuit's record also did not improve held-out predictions, which is unsurprising with fewer than ten attempts at most circuits.

In 2026 fewer undercuts worked (36\% against 50\%; Figure~\ref{fig:2026}) and leaders responded faster, with 14 of the 22 responses on the next lap; only 3 of those undercuts worked. Even with the delay and the other factors held fixed, 2026 attempts had lower odds (odds ratio 0.34), and the 2023--25 model expected 56\% success against the 36\% observed. The new tyres may take longer to warm up, which would hurt the chaser's out-lap, and teams may also have learned to respond sooner. The calibrated rule and the gap-and-delay model still beat the base rate in 2026 (Brier scores 0.201 and 0.212, against 0.250), so the arithmetic seems to have survived the new rules. We cannot yet say the odds have changed. The interval for the change in success rate runs from $-36$ to $+3$ points, which includes no change, and the 2026 indicator reaches only $p=0.049$ in a test we planned as exploratory.
.
% Edit freely; rebuild with: tectonic -X compile main.tex

The yes/no rule is overconfident mainly because it assumes that every lap before the leader responds earns the full fresh-tyre advantage, taken as a median over other cars rather than what this chaser gets on this compound. The first of those laps is the least reliable, since the chaser starts it on cold tyres while the leader can push on its in-lap, and we suspect that a good share of undercuts are decided there. Our data cannot settle whether that lap should count at all. With it, the rule's 50/50 point is a margin of 1.0~s; without it, the point moves to $-0.3$~s.

With two fitted numbers, the calibrated rule predicted as well as the full model with ten (Brier scores 0.211 and 0.218; Table~\ref{tab:rq2}), and the gap and the response delay alone did at least as well as either (0.200; Table~\ref{tab:rq5}). Since there are only 154 attempts and 77 successes, every extra coefficient carries noise, and in a new season that noise can cost more than the factor adds. The full model was the best method with 2023 held out but worse than the base rate with 2025 held out. Broadcasters could quote a probability instead of calling the outcome. Teams have more to go on (tyre temperatures, degradation models, the rival's radio and their own plans), and any of these might beat the gap alone.

Of the two numbers that matter, the chaser controls only the gap (Figure~\ref{fig:lookup}). From within a second, 79\% of undercuts worked, and the model gives better-than-even odds up to about 1.2~s even when the leader responds on the next lap. At 2~s the odds are better than even only if the leader waits three laps or more, so a chaser at the edge of the ``striking range'' is relying on a slow response. If leaders waited only when they were safe, undercuts against a slow response should rarely have worked, but 72--75\% of them did. Some leaders may be forced to wait by another car they are racing, a teammate already in the pit lane, traffic where they would rejoin or a plan for a longer stint, and others may simply misjudge the threat. Because the leader picks its delay after seeing the chaser stop, we cannot say that a late response causes a failed defence. Game-theoretic models of pit stops \citep{aguad2024pit} treat the undercut as a contest between two drivers' decisions, and our results are consistent with that.

About half of the difference between successful and failed undercuts came from the stops themselves (3.6 of 7.5~s; Figure~\ref{fig:mech}). When the leader's stop was more than 3~s slower, 15 of 17 undercuts worked. A team's pit-crew record did not predict its stops in a given attempt (correlation 0.07), so slow stops look like one-off events. An undercut that failed on a slow wheel change may therefore have been the right call when it was made, and should be judged that way. Race simulations should also treat stop time as random, as Monte Carlo simulations do \citep{heilmeier2020monte}, since fixed stop times would miss an effect of this size.

Faster cars won more undercuts (58\% against 44\%), but they also started closer, and once the gap and the response delay were in the model, pace was no longer associated with success (odds ratio 1.00 per 0.1~s a lap). The team effect that did show up was on the defending side (Figure~\ref{fig:teams}), though with fewer than 25 attempts per team and six tests added after the main analysis it is weak evidence. Undercuts against Aston Martin, Mercedes and Ferrari rarely worked, and those against Alpine, Racing Bulls and Sauber/Audi usually did. Some teams may be better at parts of a defence that timing data cannot see, such as warming their own new tyres on the out-lap. Front-runners also did worse as chasers. With the gap and the response delay held fixed, a chaser in the top six had lower odds of success (odds ratio 0.41, $p=0.06$, in groups chosen after seeing the data). The leading teams may watch each rival more closely and respond in ways the delay does not capture.

Circuits mattered mainly through the gap, with undercuts working more often where attempts started closer ($\rho=-0.54$) and no more often where tyre wear or the fresh-tyre gain was larger. A circuit's record also did not improve held-out predictions, which is unsurprising with fewer than ten attempts at most circuits.

In 2026 fewer undercuts worked (36\% against 50\%; Figure~\ref{fig:2026}) and leaders responded faster, with 14 of the 22 responses on the next lap; only 3 of those undercuts worked. Even with the delay and the other factors held fixed, 2026 attempts had lower odds (odds ratio 0.34), and the 2023--25 model expected 56\% success against the 36\% observed. The new tyres may take longer to warm up, which would hurt the chaser's out-lap, and teams may also have learned to respond sooner. The calibrated rule and the gap-and-delay model still beat the base rate in 2026 (Brier scores 0.201 and 0.212, against 0.250), so the arithmetic seems to have survived the new rules. We cannot yet say the odds have changed. The interval for the change in success rate runs from $-36$ to $+3$ points, which includes no change, and the 2026 indicator reaches only $p=0.049$ in a test we planned as exploratory.
.
% Edit freely; rebuild with: tectonic -X compile main.tex

The yes/no rule is overconfident mainly because it assumes that every lap before the leader responds earns the full fresh-tyre advantage, taken as a median over other cars rather than what this chaser gets on this compound. The first of those laps is the least reliable, since the chaser starts it on cold tyres while the leader can push on its in-lap, and we suspect that a good share of undercuts are decided there. Our data cannot settle whether that lap should count at all. With it, the rule's 50/50 point is a margin of 1.0~s; without it, the point moves to $-0.3$~s.

With two fitted numbers, the calibrated rule predicted as well as the full model with ten (Brier scores 0.211 and 0.218; Table~\ref{tab:rq2}), and the gap and the response delay alone did at least as well as either (0.200; Table~\ref{tab:rq5}). Since there are only 154 attempts and 77 successes, every extra coefficient carries noise, and in a new season that noise can cost more than the factor adds. The full model was the best method with 2023 held out but worse than the base rate with 2025 held out. Broadcasters could quote a probability instead of calling the outcome. Teams have more to go on (tyre temperatures, degradation models, the rival's radio and their own plans), and any of these might beat the gap alone.

Of the two numbers that matter, the chaser controls only the gap (Figure~\ref{fig:lookup}). From within a second, 79\% of undercuts worked, and the model gives better-than-even odds up to about 1.2~s even when the leader responds on the next lap. At 2~s the odds are better than even only if the leader waits three laps or more, so a chaser at the edge of the ``striking range'' is relying on a slow response. If leaders waited only when they were safe, undercuts against a slow response should rarely have worked, but 72--75\% of them did. Some leaders may be forced to wait by another car they are racing, a teammate already in the pit lane, traffic where they would rejoin or a plan for a longer stint, and others may simply misjudge the threat. Because the leader picks its delay after seeing the chaser stop, we cannot say that a late response causes a failed defence. Game-theoretic models of pit stops \citep{aguad2024pit} treat the undercut as a contest between two drivers' decisions, and our results are consistent with that.

About half of the difference between successful and failed undercuts came from the stops themselves (3.6 of 7.5~s; Figure~\ref{fig:mech}). When the leader's stop was more than 3~s slower, 15 of 17 undercuts worked. A team's pit-crew record did not predict its stops in a given attempt (correlation 0.07), so slow stops look like one-off events. An undercut that failed on a slow wheel change may therefore have been the right call when it was made, and should be judged that way. Race simulations should also treat stop time as random, as Monte Carlo simulations do \citep{heilmeier2020monte}, since fixed stop times would miss an effect of this size.

Faster cars won more undercuts (58\% against 44\%), but they also started closer, and once the gap and the response delay were in the model, pace was no longer associated with success (odds ratio 1.00 per 0.1~s a lap). The team effect that did show up was on the defending side (Figure~\ref{fig:teams}), though with fewer than 25 attempts per team and six tests added after the main analysis it is weak evidence. Undercuts against Aston Martin, Mercedes and Ferrari rarely worked, and those against Alpine, Racing Bulls and Sauber/Audi usually did. Some teams may be better at parts of a defence that timing data cannot see, such as warming their own new tyres on the out-lap. Front-runners also did worse as chasers. With the gap and the response delay held fixed, a chaser in the top six had lower odds of success (odds ratio 0.41, $p=0.06$, in groups chosen after seeing the data). The leading teams may watch each rival more closely and respond in ways the delay does not capture.

Circuits mattered mainly through the gap, with undercuts working more often where attempts started closer ($\rho=-0.54$) and no more often where tyre wear or the fresh-tyre gain was larger. A circuit's record also did not improve held-out predictions, which is unsurprising with fewer than ten attempts at most circuits.

In 2026 fewer undercuts worked (36\% against 50\%; Figure~\ref{fig:2026}) and leaders responded faster, with 14 of the 22 responses on the next lap; only 3 of those undercuts worked. Even with the delay and the other factors held fixed, 2026 attempts had lower odds (odds ratio 0.34), and the 2023--25 model expected 56\% success against the 36\% observed. The new tyres may take longer to warm up, which would hurt the chaser's out-lap, and teams may also have learned to respond sooner. The calibrated rule and the gap-and-delay model still beat the base rate in 2026 (Brier scores 0.201 and 0.212, against 0.250), so the arithmetic seems to have survived the new rules. We cannot yet say the odds have changed. The interval for the change in success rate runs from $-36$ to $+3$ points, which includes no change, and the 2026 indicator reaches only $p=0.049$ in a test we planned as exploratory.
